% Retained preamble for this source, sliced MECHANICALLY from its own main TeX.
% Source file : full-original-native.tex (lines 1-154 of the file)
% The block below is verbatim author text; the wrapper class, document metadata and
% \begin{document} are omitted because the assembler supplies them.
% lines taken: 1, 2, 3, 4, 5, 6, 7, 8, 9, 10, 11, ... (154 lines total)
%~Mouliné par MaN_auto v.0.25.0 2022-03-23 11:09:47
\documentclass[AHL,Unicode,longabstracts,MathTools,published]{cedram}

\usepackage{mathrsfs}
\usepackage{amscd}
\usepackage{amssymb}






\AtBeginDocument{\newtheorem{assumption}[cdrthm]{Assumption}}




\newcommand{\rme}{\mathrm{e}}
\newcommand{\ri}{\mathrm{i}}
\newcommand{\rmc}{\mathrm{c}}
\newcommand{\rms}{\mathrm{s}}

\newcommand{\Id}{\mathrm{Id}}
\newcommand{\Stokes}{\mathrm{Stokes}}
\newcommand{\Ran}{\mathrm{Ran}}
\newcommand{\divl}{\mathrm{div}}
\newcommand{\Lang}{\mathrm{Lang}}
\newcommand{\Boltz}{\mathrm{Boltz}}
\newcommand{\LSI}{\mathrm{LSI}}


\DeclareMathOperator{\Ent}{Ent}
\newcommand{\cL}{\mathcal{L}}
\newcommand{\cLs}{\mathcal{S}}
\newcommand{\cLa}{\mathcal{A}}
\newcommand{\cB}{\mathcal{B}}
\newcommand{\cX}{\mathcal{X}}
\newcommand{\cD}{\mathcal{D}}
\newcommand{\cR}{\mathcal{R}}
\newcommand{\cH}{\mathcal{H}}
\newcommand{\cLM}{\mathcal{M}}
\newcommand{\cLN}{\mathcal{N}}
\newcommand{\cLU}{\mathcal{U}}
\newcommand{\clV}{\mathcal{V}}

\newcommand{\R}{\mathbb{R}}
\newcommand{\bbZ}{\mathbb{Z}}
\newcommand{\bbT}{\mathbb{T}}
\newcommand{\bbN}{\mathbb{N}}



\newcommand{\bfun}{\mathbf{1}}
\newcommand{\scrC}{\mathscr{C}}
\newcommand{\scrT}{\mathscr{T}}

\newcommand{\frkS}{\mathfrak{S}}
\newcommand{\frka}{\mathfrak{a}}
\newcommand{\frkb}{\mathfrak{b}}

\newcommand{\Schur}{\mathfrak{S}_0}
\newcommand{\Schurb}{\mathfrak{S}_1}

\newcommand{\cLham}{\mathcal{L}_{\rm ham}}
\newcommand{\cLFD}{\mathcal{L}_{\rm FD}}

\newcommand{\eps}{\varepsilon}
\newcommand{\subplus}{\textnormal{\texttt{+}}}
\renewcommand{\leq}{\leqslant}
\renewcommand{\geq}{\geqslant}
\renewcommand{\le}{\leqslant}
\renewcommand{\ge}{\geqslant}

\newcommand{\cLNH}{\cL_{\rm NH}}
%%%%%%%%%%%%%%%%%%%%%%%%%%%%%%%%%%%%%%%%%%%%%%%%%%%%%%%%%%%
\graphicspath{{./figures/}}

\newcommand*{\mk}{\mkern -1mu}
\newcommand*{\Mk}{\mkern -2mu}
\newcommand*{\mK}{\mkern 1mu}
\newcommand*{\MK}{\mkern 2mu}

%\hypersetup{urlcolor=purple, linkcolor=blue, citecolor=red}

\newcommand*{\romanenumi}{\renewcommand*{\theenumi}{\roman{enumi}}}
\newcommand*{\Romanenumi}{\renewcommand*{\theenumi}{\Roman{enumi}}}
\newcommand*{\alphenumi}{\renewcommand*{\theenumi}{\alph{enumi}}}
\newcommand*{\Alphenumi}{\renewcommand*{\theenumi}{\Alph{enumi}}}
\let\oldtilde\tilde
\renewcommand*{\tilde}[1]{\mathchoice{\widetilde{#1}}{\widetilde{#1}}{\oldtilde{#1}}{\oldtilde{#1}}}
\let\oldforall\forall
\renewcommand*{\forall}{\mathrel{\oldforall}}
%%%%%%%%%%%%%%%%%%%%%%%%%%%%%%%%%%%%%%%%%%%%%%%%%%%%%%%%%%%%%%

\title[Hypocoercivity with Schur complements]{Hypocoercivity with Schur complements}
\alttitle{Hypocoercivité par compléments de Schur}

\subjclass{35Q84, 35H10, 35K65, 82C31, 82C40}
\keywords{Hypocoercivity, Langevin dynamics, Schur complements, resolvent estimates}


\author[\initial{E.} \lastname{Bernard}]{\firstname{\'Etienne} \lastname{Bernard}}
\address{CERMICS, \\
\'Ecole des Ponts,\\
Marne-la-Vallée (France)}
\email{etienne.bernard@enpc.fr}

\thanks{The work of A.~Levitt and G.~Stoltz was funded in part from the European Research Council (ERC) under the European Union's Horizon 2020 research and innovation programme (grant agreement No~810367). The work of G.~Stoltz was also funded by the Agence Nationale de la Recherche, under grant ANR-19-CE40-0010-01 (QuAMProcs). The work of M.~Fathi was funded by the ANR, under grants ANR-17-CE40-0030 (EFI) and ANR-18-CE40-006 (MESA), as well as ANR-11-LABX-0040-CIMI within the program ANR-11-IDEX-0002-02.}


\author[\initial{M.} \lastname{Fathi}]{\firstname{Max} \lastname{Fathi}}
\address{Université Paris Cité\\
and Sorbonne Université, CNRS,\\
Laboratoire Jacques-Louis Lions\\
and Laboratoire de Probabilités,\\
Statistique et Modélisation,\\
F-75013 Paris  (France)}
\email{mfathi@lpsm.paris}

\author[\initial{A.} \lastname{Levitt}]{\firstname{Antoine} \lastname{Levitt}}
\address{MATHERIALS project-team,\\
Inria Paris (France) \\
\& CERMICS, \\
\'Ecole des Ponts,\\
Marne-la-Vallée (France)}
\email{antoine.levitt@inria.fr}

\author[\initial{G.} \lastname{Stoltz}]{\firstname{Gabriel} \lastname{Stoltz}}
\address{CERMICS,\\
\'Ecole des Ponts,\\
Marne-la-Vallée (France)\\
\& MATHERIALS project-team, \\
Inria Paris (France)}
\email{gabriel.stoltz@enpc.fr}



\begin{abstract}
We propose an approach to obtaining explicit estimates on the resolvent of hypocoercive operators by using Schur complements, rather than from an exponential decay of the evolution semigroup combined with a time integral. We present applications to Langevin-like dynamics and Fokker--Planck equations, as well as the linear Boltzmann equation (which is also the generator of randomized Hybrid Monte Carlo in molecular dynamics). In particular, we make precise the dependence of the resolvent bounds on the parameters of the dynamics and on the dimension. We also highlight the relationship of our method with other hypocoercive approaches.
\end{abstract}

\begin{altabstract}
Nous proposons une approche fondée sur des compléments de Schur, qui permet d'obtenir des estimées explicites sur la résolvante d'opérateurs hypocoercifs, sans passer par des estimées de décroissance exponentielle du semigroupe d'évolution intégrées en temps. Nous présentons une application de cette méthode aux opérateurs de Fokker--Planck associés à des dynamiques de type Langevin, ou à l'opérateur d'évolution de l'équation de Boltzmann linéaire (qui correspond également au générateur d'un processus stochastique connu sous le nom de méthode de Monte Carlo hybride randomisée en dynamique moléculaire). Nous explicitons en particulier la dépendance des bornes de la résolvante en fonction des paramètres de la dynamique et de la dimension de l'espace ambiant. Nous discutons également les relations entre notre approche et les autres méthodes hypocoercives.
\end{altabstract}

\datereceived{2021-01-25}
\daterevised{2021-09-10}
\dateaccepted{2021-11-16}

\editor{N. Raymond}


%%%%%%%%%%%%%%%%%%%%%%%%%%%%%%%%%%%%%%%%%%%%%%%%%%%%%%%
\dateposted{2022-05-04}
